\documentclass[11pt,a4paper]{article}

\usepackage[margin=1in]{geometry}
\usepackage[T1]{fontenc}
\usepackage{lmodern}
\usepackage{microtype}
\usepackage{enumitem}
\usepackage{xurl}
\usepackage[colorlinks=true,linkcolor=black,urlcolor=blue,citecolor=blue]{hyperref}
\usepackage{titlesec}

\titleformat{\section}{\large\bfseries}{\thesection.}{0.6em}{}
\setlist{itemsep=3pt,parsep=0pt,topsep=4pt}

\title{\textbf{Weekly Progress Report}\\[0.3em]
\large Submission~1}
\author{Ahmed Khader}
\date{22 July 2026}

\begin{document}

\maketitle

\begin{abstract}
\noindent This week was devoted to building the theoretical foundation for the
project and surveying the relevant literature. The applied target is chaotic
time-series prediction, with electrocardiogram (ECG) signal forecasting as the
motivating application. The work moved from the fundamentals of chaos theory
and its predictability limits, through machine-learning approaches to
forecasting chaotic systems (with a focus on reservoir computing), to
evaluation methodology and an initial research direction. It closed with
targeted literature gathering on data-analysis pipelines and on successful
forecasting methods in adjacent fields such as numerical weather prediction.
\end{abstract}

\section{Theoretical Foundations}

\begin{itemize}
  \item Conducted a literature review framing the problem statement, and built
        up an understanding of the core theory: chaos theory, dynamical
        systems, and the Lyapunov exponent as a measure of sensitivity to
        initial conditions.
  \item Studied the central challenge of chaos theory, namely the exponential
        divergence of nearby trajectories, which imposes a fundamental
        predictability horizon, and how machine learning can be brought to
        bear on it. Related this to the target application of ECG signal
        forecasting, where the recorded dynamics are quasi-periodic and can
        exhibit chaotic behaviour.
\end{itemize}

\section{Machine Learning for Chaotic Systems}

\begin{itemize}
  \item Examined the limitations of current machine-learning approaches when
        applied to chaotic systems, and the reasons these limitations arise.
  \item Studied reservoir computing (echo state networks) as a leading
        approach for learning chaotic dynamics, and reviewed the primary
        literature on its limitations.
  \item Surveyed how the performance of machine-learning methods on chaotic
        systems is evaluated, identifying the following regimes:
        \begin{itemize}
          \item \textbf{Immediate (one-step) prediction:} forecasting the next
                value of the signal. This is the simplest regime, though it is
                not error-free: the small errors made here are what accumulate
                during autoregressive rollout and drive the failure of
                long-horizon prediction.
          \item \textbf{Valid prediction time:} how long the forecast stays
                close to the true trajectory before diverging, typically
                measured in Lyapunov times. This is the most common single
                metric in the reservoir-computing literature and bridges
                short-term accuracy and long-term behaviour.
          \item \textbf{Long-term qualitative agreement:} whether the forecast
                stays on the correct attractor and reproduces the overall
                pattern of the dynamics, rather than the exact values. For ECG,
                this is the clinically meaningful question, such as whether the
                forecast trends in the same direction as a patient entering a
                cardiac event.
          \item \textbf{Long-term pointwise accuracy:} exact value-by-value
                agreement over a long horizon. This is extremely difficult
                because of sensitivity to initial conditions: a positive
                Lyapunov exponent causes small errors to grow exponentially,
                which sets a hard predictability limit.
        \end{itemize}
\end{itemize}

\section{Initial Research Direction}

\begin{itemize}
  \item Formed an initial hypothesis on how the underlying prediction problem
        might be improved, and on the potential role of the effective
        (finite-time) Lyapunov exponent in characterising and improving
        forecast performance. This remains an early idea to be developed
        further.
\end{itemize}

\section{Literature Gathering}

\begin{itemize}
  \item Collected papers on data analysis in order to design robust and
        efficient pipelines for processing the project data.
  \item Collected papers directly addressing the target problem and the ways
        the existing literature has approached it.
  \item Began surveying adjacent fields, notably numerical weather prediction,
        to identify transferable methods. Successful examples reviewed include
        graph neural network models (such as GraphCast) and 3D transformer
        architectures (such as Pangu-Weather's 3D Earth-Specific Transformer).
\end{itemize}

\section{Planned for Next Week}

\begin{itemize}
  \item Read and summarise the gathered papers in depth, with written feedback
        on each.
  \item Develop the initial research direction into a more concrete proposal.
\end{itemize}

\end{document}
